\section{Cyclic sequences}
\todo{Lecture 7}
\begin{definition}
Let $A$ be a set. A linear sequence of length $n$ in alphabet $A$ is an element of $A^{n}$ :
$$
a=(a_{1},\dots,a_{n}),\quad a_{k}\in A,~k=1,\dots,n
$$
The number of linear sequences of length $n$ in an alphabet of size $r$ is $r^{n}$.

Consider the following equivalence relation on the set of linear sequences 
$$
(a_{1},\dots,a_{n})\sim (a_{2},a_{3},\dots,a_{n},a_{1})
$$
Two linear sequences are equivalent if one of them can be obtained from another by a cyclic shift.
\end{definition}
\begin{reminder}[Equivalence relation]
Let $X$ be a set. A binary relation on $X$ is a subset of $X\times X$. A relation $R\subset X$ is an equivalent relation if it is:
\begin{enumerate}
\item Reflexive: $\forall x\in X: (x,x)\in R$.
\item Symmetric: $\forall x,y\in X:(x,y) \in R \Rightarrow(y,x)\in R$.
\item Transitive: $\forall x,y,z\in Z:(x,y)\in R,(y,z)\in R\Rightarrow(x,z)\in R$.
\end{enumerate}
\end{reminder}
\begin{definition}
A cyclic sequence of length $n$ in alphabet $A$ is an equivalence class of linear sequences with respect to the relation $\sim$.
\end{definition}
\begin{example}
There are 8 linear sequences of length 3 in alphabet $\{ a,b \}$ and only 4 cyclic sequences.
\end{example}
\begin{proposition}
The number $T(n,r)$ of cyclic sequences of length $n$ on an alphabet of size $r$ is
$$
T(n,r)=\frac{1}{n}\sum _{d\mid n} \phi(n/d)r^{d}
$$
where $\phi$ is the Euler totient function.
\end{proposition}
\begin{proof}
\begin{definition}
A period of a cyclic sequences $(a_{1},\dots,a_{n})$ is a minimal number $k\in \{ 1,\dots,n \}$ such that
$$
(a_{1},\dots,a_{n})=(a_{1+k},\dots,a_{1},\dots,a_{k})
$$
are equal as linear sequences.
\end{definition}
Let $M(d,r)$ be the number of cyclic sequences of length $d$ and period exactly $d$. The following identity holds
$$
r^{n}=\sum _{d\mid n}d M(d,r).
$$
We have the map
$$
\{ \text{Linear sequences of length }n\text{ in alphabet } [r]  \}\overset{\pi}\to \{ \text{Cyclic sequences of length }n\text{ in } [r] \}
$$
where $\pi$ is the projection into an equivalency class. The number of preimages of a cyclic sequences under the map $\pi$ is $d$, the period of the sequence.

Therefore
$$
|L(n,r)|=\sum _{d\mid n}d|M(d,r)|.
$$


Starting with the identity 
$$
r^{n}=\sum _{d\mid n} dM(d,r)
$$
we apply the Mobius inversion formula and obtain
\begin{equation}
nM(n,r)=\sum _{d\mid n}\mu(d/n)r^{d}. \label{mobinv}
\end{equation}
Each cyclic sequence has a well defined period $d$ and it corresponds to the unique cyclic sequence of length $d$ and period $d$. Thus
\begin{equation}
T(n,r)=\sum _{d\mid n}M(d,r).n \label{mobinvtwo}
\end{equation}
Now we substitute \ref{mobinv} into \ref{mobinvtwo}:
\begin{align*}
T(n,r) & =\sum _{d\mid n}M(d,r) \\
 & =\sum _{d\mid n} \frac{1}{d}\sum _{d'\mid d}\mu(d/d')r^{d'} \\
 & =\sum _{d'\mid n}\left( \sum _{d''\mid (n/d')} \frac{1}{d'd''} \mu(d'')  \right)r^{d'}
\end{align*}
Now it remains to compute the sum
$$
\sum _{d''\mid (n/d')} \frac{1}{d'd''} \mu(d'')
$$
\begin{exercise}
Show that for $n\in \mathbf{Z}_{\ge 1}$, $\sum _{d\mid n} \frac{1}{d}(\mu(d))=\frac{\phi(n)}{n}$.
\end{exercise}
Using this exercise we compute
$$
T(n,r)=\sum _{d'\mid n} \frac{\phi(n/d')}{n}r^{d'}
$$
\end{proof}

\section{Partially ordered set (posets)}
\begin{definition}
A binary relation on a set $A$ is a subset $R\subseteq A\times A$.
\end{definition}
\begin{definition}
A relation is antisymmetric provided $(a,b)\in R$ and $(b,a)\in R$ imply $a=b$.
\end{definition}
\begin{definition}
A relation is transitive if $(a,b)\in R$ and $(b,c)\in R$ imply $(a,c)\in R$.
\end{definition}
\begin{definition}
A relation is reflexive if $(a,a)\in R$ for all $a\in R$.
\end{definition}
\begin{definition}
A partial order on a set $A$ is an antisymmetric, reflexive and transitive relation $R\subseteq A\times A$. 
\end{definition}
\begin{definition}
A partially ordered set (or poset for short) is a set together with a partial order.
\end{definition}
\begin{example}
Let $A$ be a set. The set $2^{A}$ is partially ordered by inclusion $\subseteq$.
\end{example}
\begin{example}
The set $\mathbf{Z}_{\ge 1}$ is partially ordered by the relation "$d$ divides $n$".
\end{example}
\begin{definition}
A partially ordered set $(X,\le)$ is locally finite if for all $x,y\in X$ the interval
$$
[x,y]:= \{ z\in X:y\le z \le x \}
$$
is finite.

We say that $0\in X$ is a zero element if $0\le x$ for all $x \in X$.
\end{definition}

\subsection{Mobius inversion formula for posets}

Let $(X,\le )$ be a partially ordered locally finite set with 0. Suppose that $f:X\to \mathbf{C}$ is a function. We define a new function $F:X\to \mathbf{C}$ by
$$
F(x):= \sum _{y\le x}f(y)
$$
How do recover $f$ from $F$?
\begin{theorem}[Mobius inversion for posets] \label{mobposet}
Given a partially ordered set $X$, there is a two-variable function $M:X\times X\to \mathbf{R}$ such that
$$
F(x)=\sum _{y\le x}f(y)\Leftrightarrow f(x)=\sum _{y\le x}F(y)M(y,x)
$$
$M$ is called the Mobius function of the poset.
\end{theorem}
\begin{definition}
Given a partially ordered set $X$, the incidence algebra $A(X)$ is the set of complex-valued functions $f:X^{2}\to \mathbf{C}$ satisfying $f(x,y)=0$ unless $x\le y$. 
\end{definition}
$A(X)$ is a vector space over $\mathbf{C}$ with respect to pointwise addition and multiplicavtion by scalars. To make it into an "algebra" we need one more operation:
\begin{definition}
Given $f,g\in A(X)$, their convolution $f*g$ is defined by
$$
f*g(x,y)=\sum _{z:x\le z \le y} f(x,z)g(z,y)
$$
\end{definition}
\begin{definition}
The delta function $\delta$ is the following element of $A(X)$:
$$
\delta(x,y)=\begin{cases}
1 & \text{if }x=y  \\
0 & \text{otherwise}.
\end{cases}
$$
\end{definition}
\begin{lemma}
A function $f\in A(X)$ has a left and right inverse with respect to the convolution if and ony if $f(x,x)\neq0$ for all $x \in X$.
\end{lemma}
\begin{proof}
Given $f\in A(X)$ we find $g\in A(X)$ such that
$$
f*g(x,y)=\sum _{x\le z \le y}f(x,z)g(z,y)=\delta(x,y)
$$
For all $x \in X$ we have that $f*g(x,x)=\delta(x,x)=1$. Therefore, the condition $f(x,x)\neq0$ for all $x \in X$ is necessary for the existence of the inverse.

We define $g(x,x):=f(x,x)^{-1}$ for all $x \in X$. To define $g(x,y)$ for $x<y$, we assume by induction, that we have already found all $g(z,y)$ for all $z$ satisfying $x<z\le y$. Then
$$
\delta(x,y)=0=\sum _{x\le z \le y}f(x,z)g(z,y)-f(x,x)g(x,y)=\sum _{x<z \le y}f(x,z)g(z,y)
$$
We can find $g(x,y)$ from this identity because $f(x,x)\neq 0$ and we know all the terms of the finit sum on the right hand side.

Finally, if $f*g_{1}=\delta$ and $g_{2}*f=\delta$ then
$$
g_{2}=g_{2}*\delta=g_{2}*f*g_{1}=(g_{2}*f)*g_{1}=\delta*g_{1}=g_{1}.
$$
Therefore, the left and the right inverses of $f$ coincide.
\end{proof}
\begin{definition}
The zeta function $Z(x,y)$ of the poset $(X,\le)$ is the function
$$
Z(x,y)=\begin{cases}
1 & \text{if }x\le y \\
0 & \text{otherwise}.
\end{cases}
$$

The Mobius function $M(x,y)$ is the inverse of the zeta function $Z$ with respect to convolution.
\end{definition}
\begin{proof}[Proof of theorem \ref{mobposet}]
Let $f:X\to \mathbf{C}$ be a function and $F:X\to \mathbf{C}$ is defined by 
$$
F(x)=\sum _{y\le x}f(y)
$$
Then for a fixed $x \in X$:
\begin{align*}
\sum _{y\le x}F(y)M(y,x) & =\sum _{y\le x}M(y,x)\sum _{z\le y}f(z) \\
 & =\sum _{z\le y \le x}f(z)M(y,x) \\
 & =\sum _{z\le x}f(z)\sum _{y:z\le y \le x}M(y,x) \\
 & =\sum _{z \le x}f(z)\left( \sum _{z\le y\le x} Z(z,y)M(y,x) \right) \\
 & =\sum _{z\le x}f(z)(Z*M)(z,x) \\
 & =\sum _{z\le x}f(z)\delta(z,x)=f(x)
\end{align*}

Now we show that the inverse is also true. Let $F:X\to \mathbf{C}$ be a function and we define $f:X\to \mathbf{C}$ by
$$
f(x)=\sum _{y\le x} F(y)M(y,x)
$$
foral $x \in X$.

Then
\begin{align*}
\sum _{y\le x} & =\sum _{y\le x}\sum _{z \le y}F(z)M(z,y) \\
 & =\sum _{z,y:z\le y\le x}F(z)M(z,y) \\
 & =\sum _{z\le y\le x}F(z)M(z,y)Z(y,x) \\
 & =\sum _{z\le x}F(z)\left( \sum _{z\le y \le x}M(z,y)Z(y,x) \right) \\
 & =\sum _{z\le x}F(z)\delta(z,x)=F(x)
\end{align*}
\end{proof}

\begin{lemma}
Let $2^{A}$ be the set of subsets of a finite set $A$. $2^{A}$ is partially ordered by inclusion. The Mobius function of $2^{A}$ is given by
$$
M(x,y)=(-1)^{|x|-|y|}
$$
where $x\subseteq y\subseteq A$.
\end{lemma}
\begin{proof}
We have to show that $M*Z=\delta$. Let $x,y\subseteq A$ such that $x\subseteq y$. We compute
\begin{align*}
M*Z(x,y)=\sum _{z:x \subseteq z \subseteq y} M(x,z)Z(z,y) \\
 & =\sum _{x \subseteq z \subseteq y}(-1)^{|x|-|z| } \\
 & =\sum _{w\subseteq y\setminus x} (-1)^{|w|}=\begin{cases}
0 & \text{if }|y\setminus x| \ge 1 \\
1 & \text{if }|y\setminus x|=\varnothing.
\end{cases}
\end{align*}
\end{proof}
\begin{lemma}
Consider the set $\mathbf{Z}_{\ge 1}$, partially ordered by the relation "$x$ divides $y$".

then the Mobius function of this poset is:
$$
M(x,y)=\mu(y/x).
$$
\end{lemma}
\begin{proof}
We need to show that $M*Z=\delta$.

For $x,y\in \mathbf{Z}_{\ge 1}$ with $x\mid y$ we have
\begin{align*}
M*Z(x,y) & =\sum _{\underset{x\le z \le y}{z\in \mathbf{Z}_{\ge 1}}} M(x,z)Z(z,y)  \\
 & =\sum _{z:x \mid z, z\mid y} \mu(z/x) \\
 & =\sum _{d\mid (y/x)} \mu(d)=
\begin{cases}
1 & \text{if } \frac{y}{x}=1 \\
0 & \text{otherwise.}
\end{cases}
\end{align*}
\end{proof}